\chapter{总结}
\label{chap:conclusion}

\section{研究结论}
\label{sec:fivepone}
将高能对撞机环境下产生的粒子对的自旋视为二量子比特系统，对大量事件统计所得到的混合态内部可能存在量子纠缠。本研究利用基于扩散模型和Transformer模型构建的OmniLearn生成模型，对$pp \to \tau^+ \tau^- \to (\pi^+ \bar{\nu}_{\tau}) (\pi^- \nu_{\tau})$ 过程产生的中微子的动量信息进行重建，其后应用量子层析方法确定$\tau^+ \tau^-$ 粒子对的各自旋关联矩阵元，最终通过纠缠判据$\mathcal{O}_{E}^{\pm}=0.1688>0 $确定该生成模型所得的$\tau^+ \tau^- $对间存在量子纠缠。此外，本研究将所使用的生成模型与传统的MMC方法所得结果进行比较，进一步印证了所使用的机器学习方法在粒子物理实验海量数据处理上的可靠性。

\section{研究意义和价值}
\label{sec:fiveptwo}
高能对撞条件下的体系的纠缠为研究者提供了新的视角。本研究关注并尝试确认$\tau^+ \tau^-$ 粒子对之间的纠缠，这对于量子力学基本问题的回答具有重要意义，并启发后人进一步探索将量子纠缠作为手段和工具应用在寻找新物理上的可能性。

\section{展望和下一步工作计划}
\label{sec:fivepthree}
在研究过程中各粒子四动量生成结果往往与蒙特卡洛模拟中的重建层级数据相比较，但实际上将其与真实层级数据相比较更能检验生成结果的合格与否。这要求对机器学习结果进行解卷积以去除探测器效应带来的影响，还原出真实层级上的分布。

此外，本研究对机器学习结果的不确定度分析是粗糙的，更合理的做法是通过诸如误差传递等方式考虑$p_T$、$\phi$ 、$\theta$等参量的误差如何影响自旋关联矩阵，进而影响纠缠判断的有效性。以上两点将作为未来工作的重点方向持续推进，以期取得更加卓越的成果。